% Ajustes de apresentação para o PDF acadêmico em ABNT.
\usepackage{indentfirst}
\setlength{\parindent}{1.25cm}
\setlength{\parskip}{0pt}

% Tipografia pensada para leitura em tela: o livro é distribuído em PDF e HTML,
% não impresso. XCharter (extensão do Bitstream Charter, de Matthew Carter)
% tem hastes robustas e contraste baixo, desenhados para sobreviver a
% rasterização grosseira — o oposto do Latin Modern, herdeiro do Computer
% Modern, cujos traços finos afinam demais em tela. Fira Mono, mais escuro e
% aberto que o mono padrão, equilibra o peso do texto nos exemplos de código.
\setmainfont{XCharter}[
  Extension      = .otf,
  UprightFont    = *-Roman,
  ItalicFont     = *-Italic,
  BoldFont       = *-Bold,
  BoldItalicFont = *-BoldItalic,
]
\setmonofont{FiraMono}[
  Extension   = .otf,
  UprightFont = *-Regular,
  BoldFont    = *-Bold,
  Scale       = 0.86,
]

% Títulos com a mesma serifada do corpo, em vez do sem-serifa padrão da classe.
\setkomafont{disposition}{\rmfamily\bfseries}

% Tabelas: o espaçamento de 1,5 do corpo do texto deixa as células frouxas.
% Dentro das tabelas usa-se entrelinha simples e ganha-se legibilidade com
% uma folga maior entre as linhas, não dentro delas.
\usepackage{etoolbox}
\AtBeginEnvironment{longtable}{%
  \setstretch{1.0}%
  \setlength{\extrarowheight}{6pt}%
}
\renewcommand{\arraystretch}{1.2}

% Blocos de código também pedem entrelinha simples: a indentação já organiza a
% leitura, e o espaçamento do texto corrido só afasta linhas que devem ser
% lidas como um bloco.
\AtBeginEnvironment{Shaded}{\setstretch{1.0}}

% NBR 14724: numeração em algarismos arábicos, no canto superior direito da
% folha, a partir da primeira folha da parte textual (a capa não é numerada).
% Cabeçalho corrente com o capítulo atual nas páginas de continuação; nas
% páginas de abertura de capítulo (estilo "plain"), mostra-se só o número.
\usepackage{fancyhdr}
\pagestyle{fancy}
\fancyhf{}
\fancyhead[L]{\nouppercase{\leftmark}}
\fancyhead[R]{\thepage}
\renewcommand{\headrulewidth}{0pt}
\renewcommand{\footrulewidth}{0pt}
\fancypagestyle{plain}{%
  \fancyhf{}
  \fancyhead[R]{\thepage}
  \renewcommand{\headrulewidth}{0pt}
}

% NBR 14724: folha de rosto com instituição, curso, natureza do trabalho,
% cidade e ano, além de título e autoria.
\makeatletter
\renewcommand{\maketitle}{%
  \begin{titlepage}
    \centering
    \vspace*{1.5cm}
    {\large UNIVERSIDADE FEDERAL DO TOCANTINS\par}
    \vspace{0.15cm}
    {\large Câmpus de Palmas\par}
    \vspace{0.15cm}
    {\large Curso de Ciência da Computação\par}
    \vfill
    {\Huge\bfseries \@title \par}
    \vspace{1.2cm}
    {\large \@author \par}
    \vfill
    {\normalsize Material didático da disciplina Introdução à Ciência da\\
      Computação, do Curso de Ciência da Computação da Universidade\\
      Federal do Tocantins (UFT), Câmpus Palmas.\par}
    \vfill
    {\large Palmas -- TO\par}
    {\large \@date \par}
  \end{titlepage}
}
\makeatother
